\section{Conclusion}
\label{sec:conclusion}
Finding where a shortcut acts and knowing how to remove it turned out to be different questions. The tools answered the
first one consistently: in our models the name takes over in the last three blocks. But that answer did not settle the
second. Whether retraining succeeds depended on the repair data, which has to contain examples that contradict the
shortcut; given those, any block can learn to cancel it, and without them none does. Location mattered only for what
happens before the repair is complete: how quickly a block gets there and by which route. Even there the tools told half
the story, pointing to the late blocks and missing block~0, which repairs the shortcut where the name enters.

For using interpretability to guide repairs, this suggests a division of labour. The repair data should be checked
first, because no choice of block makes up for missing counter-examples. Localization helps most where no training is
involved, as with the read-out taken just before the located blocks; under a small training budget it points to some of
the fast blocks, but not all. It also has a clear limit: for memorized wrong labels, which share no cue, the block does
matter, and the tools did not find it.

These conclusions rest on GPT-2 small and one shortcut that we inserted ourselves, with exploratory checks on a second
cue and a second dataset. The next steps are a blind test in a second model, where the fast blocks and the read-out
layer are predicted from the interpretability tools before any repair is run, and the same questions for shortcuts that
arise naturally.
